\documentclass[10pt]{article}

\usepackage[letterpaper, left=0.6in, right=0.6in, bottom=0.4in, top=0.4in]{geometry}
\usepackage{fancyhdr}
\usepackage{titling}
\usepackage{titlesec}
\usepackage{hyperref}
\usepackage{xcolor}
\usepackage{lastpage}
%\usepackage{lipsum}
\usepackage[fleqn]{amsmath}
\usepackage{setspace}
\usepackage{stmaryrd}
\usepackage{expex}
\usepackage{tikz}
\usepackage{hanging}

%\usepackage[hmargin=1in,vmargin=0.75in]{geometry}
\usepackage{enumitem}

\usepackage{tabularx}
\usepackage{multirow}

%%% for math
\usepackage{mathtools, amssymb, stmaryrd}
\usepackage{delimset}


%\renewcommand{\familydefault}{\sfdefault}

\usepackage{libertine}
\usepackage[libertine]{newtxmath}

\usepackage{bold-extra}

\titleformat*{\section}{\Large\scshape}
\titleformat*{\subsection}{\large\scshape}
%\titleformat*{\subsubsection}{\small\bfseries}

\usepackage{framed}
\usepackage{colortbl}

\definecolor{bluegray}{rgb}{0.4, 0.6, 0.8}

\renewenvironment{leftbar}[1][\hsize]
{%
    \def\FrameCommand
    {%
        {\color{pink}\vrule width 3pt}%
        \hspace{1em}%must no space.
    }%
    \MakeFramed{\hsize#1\advance\hsize-\width\FrameRestore}%
}
{\endMakeFramed}

\hypersetup{
    colorlinks,
    linkcolor={gray},
    citecolor={gray},
    urlcolor={gray}
}

\begin{document}

\thispagestyle{empty}

\section*{$\blacksquare$ \hspace{1em} Ian Thomas White \hspace{1em} $\circ$ \hspace{1em} Computational Linguist \& Writer}

\vspace{-1em}

\leftbar[0.99\linewidth]

\begin{tabular}{lllllllll}
      \href{https://www.ianthomaswhite.com}{Website} & $\circ$ & \href{https://www.linkedin.com/in/ianthomaswhite/}{Linkedin}& $\circ$ & \href{mailto:ianthomaswhite@gmail.com}{ianthomaswhite@gmail.com} & $\circ$ & 515-508-9805 & $\circ$ & Los Angeles (open to relocation)\\
\end{tabular}

\endleftbar

\subsection*{$\square$ \hspace{1.4em} Overview}

\vspace{-1em}

\leftbar[0.99\linewidth]

\begin{tabularx}{\textwidth}{X}
    I am an academically trained linguist with substantial tech-industry experience. I have training in all linguistic subfields, with a focus on formal language theory, and I've leveraged my linguistics training at ML/AI companies. I'm a rigorous problem-solver across both quantitative (Python, SQL, Haskell) and qualitative (writing, annotation, ontology) domains. \\
\end{tabularx}

\endleftbar

\subsection*{$\square$ \hspace{1.4em} Proficiencies}

\vspace{-1em}

\leftbar[0.99\linewidth]

\begin{tabularx}{\textwidth}{X!{\color{pink}\vrule}lX!{\color{pink}\vrule}lX!{\color{pink}\vrule}lX!}
     \scshape{Programming} & & \scshape{Linguistics} & & \scshape{Languages} & & \scshape{Skills} \\
      & & & & & & \\
     \raggedright Git, Haskell, JSON, \LaTeX, Markdown, MS 365, NLTK, Notion, Python, PyTorch, Prompt Engineering, R, SQL & & \raggedright Computational, Grammar, Lexical, Morphology, NLP, Phonology, Pragmatics, Semantics, Syntax & & \raggedright French, Greek (Ancient), Latin (Classical) & & \raggedright Communication, Data Analysis \& Annotation, Ontology, Research, Statistics, Taxonomy, Technical Writing
\end{tabularx}

\endleftbar

\subsection*{$\square$ \hspace{1.4em} Experience}

\vspace{-1em}

\leftbar[0.99\linewidth]

\begin{tabularx}{\textwidth}{lX!}
    \multicolumn{2}{l}{\textsc{Linguistics Trainer and QA Specialist, Human Data} @ SpaceXAI: 08.2024-08.2026} \\
    $\circ$ & Annotated textual and visual/audio LLM outputs for adherence to phonological, syntactical, and semantic norms. \\
    $\circ$ & Conducted model evaluation and QA for Python and SQL outputs. \\
    $\circ$ & Wrote novel guidance and QA rubrics in Markdown, JSON, and Excel. \\
    $\circ$ & Collaborated with engineers to ensure labeling guidelines accorded with machine learning principles and goals.\\
    $\circ$ &  Trained new tutors and developed algorithmic processes to improve tutor workflow and inter-annotator agreement. \\
    \multicolumn{2}{l}{} \\
\end{tabularx}

\begin{tabularx}{\textwidth}{lX!}
    \multicolumn{2}{l}{\textsc{Linguistics Instructor} @ UCLA: 09.2022-12.2023} \\
    $\circ$ & Taught weekly seminars covering major linguistic subfields: phonology, morphology, syntax, semantics, pragmatics.\\
    $\circ$ & Developed unique teaching materials and technical documentation using \LaTeX\,and Haskell. \\
    $\circ$ & Evaluated student performance and conducted regular one-on-one meetings with students. \\ 
    \multicolumn{2}{l}{} \\
\end{tabularx}

\begin{tabularx}{\textwidth}{lX!}
    \multicolumn{2}{l}{\textsc{Computational Linguist} @ App Orchid: 06.2021-09.2022} \\
    $\circ$ & Designed syntactic-tree-based regexes to match complex, embedded legal phrases in an LLM pipeline. \\
    $\circ$ & Labeled and taxonomized linguistic regularities in legal corpora for POS, NER, and dependency/constituency parsing. \\
    $\circ$ & Researched string parsing techniques to optimize regex search algorithms. \\
    $\circ$ & Integrated regexes into the LLM pipeline using Python, SQL, and Git and wrote accompanying technical documentation.\\
    $\circ$ & Communicated regularly with engineers, clients, and other stakeholders to ensure optimal ML outcomes.\\
    \multicolumn{2}{l}{} \\
\end{tabularx}

\begin{tabularx}{\textwidth}{lX!}
    \multicolumn{2}{l}{\textsc{Literature Instructor} @ University of Minnesota: 09.2016-05.2020} \\
    $\circ$ & Designed and implemented syllabi for key departmental courses on literary theory. \\
    $\circ$ & Wrote from scratch unique teaching materials on grammar and rhetoric.\\
\end{tabularx}

\endleftbar

\subsection*{$\square$ \hspace{1.4em} Education}

\vspace{-1em}

\leftbar[0.99\linewidth]

\begin{tabularx}{\textwidth}{lX!}
    \multicolumn{2}{l}{\textsc{PhD (Incomplete), Linguistics} @ UCLA: 2020-2023} \\
    $\circ$ & Took graduate courses in all linguistic subfields, as well as in algorithms and formal language theory. \\
    $\circ$ & Researched and wrote on the generative capacities of Combinatory Categorial Grammars and Minimalist Grammars.\\
    $\circ$ & Applied research toward NLP/NLU applications using Python and Haskell.\\
    \multicolumn{2}{l}{} \\
\end{tabularx}

\begin{tabularx}{\textwidth}{lX!}
    \multicolumn{2}{l}{\textsc{MA, Linguistics} @ University of Minnesota: 2016-2020} \\
    $\circ$ & Thesis: ``The Argument Structure of Deverbal Nouns''. \\
    $\circ$ & Took graduate courses in all linguistic subfields: phonology, syntax, semantics, pragmatics, computational. \\
    $\circ$ & Completed corpora analysis projects using Python, NLTK, and Google Colab.\\
    \multicolumn{2}{l}{} \\
\end{tabularx}

\begin{tabularx}{\textwidth}{lX!}
    \multicolumn{2}{l}{\textsc{MA, Literature} @ University of Minnesota: 2016-2020} \\
    $\circ$ & Thesis: ``Discrete Infinity in Penelope’s Weaving Trick''. \\
    $\circ$ & Took graduate courses in literary theory, French-language poetry, and Greek-language literature.\\
    $\circ$ & Used APA, MLA, and Chicago manuals of style to write papers.\\
    \multicolumn{2}{l}{} \\
\end{tabularx}

\begin{tabularx}{\textwidth}{lX!}
    \multicolumn{2}{l}{\textsc{BA, Writing} @ Johns Hopkins University: 2008-2011} \\
    $\circ$ &  Thesis: ``Modern Mimesis: Alt-lit and Autofiction''. \\
\end{tabularx}

\endleftbar



\end{document}